\chapter{Dispersion et séries regroupées en classes}
\section*{Exercice 19.1 -- Calculer une dispersion}
$\bar x=4$. Par les écarts, $V=(4+0+0+4)/4=2$. Par la moyenne des carrés, $V=(4+16+16+36)/4-4^2=18-16=2$. Ainsi $\sigma=\sqrt2\simeq1{,}414$. L'unité de $\sigma$ est celle des données ; celle de $V$ est son carré.

\section*{Exercice 19.2 -- Changer d'unité}
La moyenne devient $1{,}8\times18+32=64{,}4$. Chaque écart à la moyenne est multiplié par $1{,}8$ : $(1{,}8C_i+32)-(1{,}8\bar C+32)=1{,}8(C_i-\bar C)$. Donc l'écart-type est $1{,}8\times3=5{,}4$. Une translation ne change pas la dispersion.

\section*{Exercice 19.3 -- Classes et histogramme}
L'effectif total est $20$. Les hauteurs peuvent être les effectifs $4,8,8$, puisque les largeurs sont égales. Le polygone joint $(0;0)$, $(10;0{,}2)$, $(20;0{,}6)$ et $(30;1)$. Les centres donnent $\bar x\simeq(4\times5+8\times15+8\times25)/20=17$. La médiane estimée est dans la deuxième classe : $10+10(0{,}5-0{,}2)/(0{,}6-0{,}2)=17{,}5$. Ces nombres sont des estimations : les effectifs par classe ne déterminent pas les valeurs individuelles.

\section*{Exercice 19.4 -- Ce que le regroupement perd}
$\bar x=(4\times4+8\times14+8\times26)/20=336/20=16{,}8$. L'écart-type reste inconnu : dans la première classe, quatre $4$ et les valeurs $2,4,4,6$ donnent la même moyenne mais des dispersions différentes. On peut garder identiques toutes les autres observations. Les moyennes de classe ne suffisent donc pas à déterminer la moyenne des carrés.

\section*{Exercice 19.5 -- Choisir les indicateurs}
Dans les deux cas la médiane vaut $12$, $Q_1$ est la valeur de rang $2$, soit $11$, et $Q_3$ celle de rang $4$, soit $13$. L'écart interquartile vaut $2$. Les moyennes sont $12$ et $146/5=29{,}2$ ; les étendues sont $4$ et $90$. La valeur $100$ modifie fortement moyenne et étendue sans modifier ici médiane et quartiles. Il faut décrire à la fois le centre, la dispersion et les valeurs extrêmes.

\section*{Exercice 19.6 -- Décomposer une variance (facultatif)}
La moyenne globale est $5$. Pour chaque observation, écrire $x-5=(x-\bar x_g)+(\bar x_g-5)$ et développer le carré. Dans chaque groupe, la somme des $x-\bar x_g$ est nulle : le terme croisé disparaît. La variance totale est donc $(1+4)/2+((2-5)^2+(8-5)^2)/2=2{,}5+9=11{,}5$. La dispersion entre les deux moyennes doit s'ajouter à la dispersion interne.
